% ============================================================
%  Chapter 21 — Multi-Qubit Entangled CIIR Controller
%               with Fractal Encoding
% ============================================================
\chapter{Multi-Qubit Entangled CIIR Controller with Fractal Encoding}
\label{ch:multi-qubit-controller}

\begin{center}
\textit{%
``Controllability is not granted by the Hamiltonian alone---it
is negotiated at the intersection of coherence time, channel
bandwidth, and error suppression depth.''}
\end{center}

% ============================================================
\section*{Abstract}
\addcontentsline{toc}{section}{Abstract}
\label{sec:ch21-abstract}

We present a complete, mathematically rigorous architecture for a
\emph{multi-qubit entangled CIIR controller} that couples a biological
interface to an $N$-qubit quantum register through a layered stack
comprising protein-based molecular qubits, a photonic/phononic coupling
network, a neural-synthetic brain--computer interface (BCI), a fractal
recursive error-correction code, and a real-time adaptive feedback
controller.  Two central theorems are established within the CIIR
categorical framework: (i) the \textbf{CIIR Controllability Theorem},
which gives necessary and sufficient conditions for universal unitary
control via Lie algebra closure; and (ii) the \textbf{CIIR Threshold
Theorem}, which proves the existence of a fault-tolerance threshold
$p_*$ below which the logical error probability decays
super-exponentially with code level $n$.  Stability is characterised by
a Lyapunov functional on the hybrid state space
$(x, \rho_Q) \in \mathbb{R}^m \times \mathcal{D}(\mathcal{H}_Q)$.
Failure modes are enumerated as strict inequality conditions, and a
four-stage experimental validation pathway is proposed.

% ============================================================
\section{Introduction and Theoretical Context}
\label{sec:ch21-intro}

The CIIR framework, developed in preceding chapters, defines a category
$\mathbf{CIIR}$ whose objects are the Hilbert spaces
$\{\mathcal{H}_Q, \mathcal{H}_B, \mathcal{H}_E\}$ and whose morphisms
are completely positive trace-preserving (CPTP) maps, control
operators, and encoding isometries.  Chapters~13--15 established quantum
mechanics as a fixed-point of the CIIR filtering functor; Chapters~16--19
extended the formalism across domains; Chapter~20 introduced the MPPK
law-discovery engine.

The present chapter addresses the \emph{engineering instantiation}
of this abstract framework.  Specifically, we ask: given the full
tripartite Hilbert space $\mathcal{H}_{\mathrm{tot}} = \mathcal{H}_Q \otimes
\mathcal{H}_B \otimes \mathcal{H}_E$, what architecture and which
parameter regimes guarantee (a) universal controllability, (b)
fault-tolerant operation, and (c) asymptotic stability under biological
noise?

The three pillars of the answer are:
\begin{enumerate}[nosep]
  \item \textbf{Entanglement-assisted control.}  Multi-qubit entanglement
        enlarges the reachable set and enables collective error detection
        that is impossible qubit-by-qubit.
  \item \textbf{Fractal error correction.}  A recursively defined code
        family $\{\mathcal{E}_n\}_{n\geq 0}$ achieves super-exponential
        logical error suppression below the threshold $p_*$.
  \item \textbf{Biological interface channel.}  The neural-to-quantum
        channel $\Phi_t$ must satisfy bandwidth and latency constraints
        that are derived from the control Lyapunov conditions.
\end{enumerate}

% ============================================================
\section{Formal Definitions}
\label{sec:ch21-defs}

All symbols are declared here and used consistently throughout.

\begin{definition}[State spaces]\label{def:state-spaces}
\begin{align*}
  \mathcal{H}_Q &\cong (\mathbb{C}^2)^{\otimes N},
    & \dim \mathcal{H}_Q &= 2^N,\\
  \mathcal{H}_B &: \text{biological interface (finite-dimensional effective space)},\\
  \mathcal{H}_E &: \text{environment (infinite-dimensional; treated via Born--Markov
                  approximation)}.
\end{align*}
The \emph{total} space is $\mathcal{H}_{\mathrm{tot}} = \mathcal{H}_Q \otimes
\mathcal{H}_B \otimes \mathcal{H}_E$.  A \emph{state} on any factor is
a positive, trace-one operator $\rho \geq 0$, $\operatorname{Tr}\rho = 1$.
\end{definition}

\begin{definition}[Protein qubit]\label{def:protein-qubit}
A \emph{protein qubit} is an effective two-level subsystem
$\mathcal{H}_{\mathrm{prot}} = \operatorname{span}\{|g\rangle, |e\rangle\}$
embedded in a larger molecular Hilbert space by the isometry
\[
  V : \mathbb{C}^2 \to \mathcal{H}_{\mathrm{prot}},
  \qquad
  V|0\rangle = |g\rangle,\quad V|1\rangle = |e\rangle,
  \qquad V^\dagger V = \mathbf{1}_2.
\]
\end{definition}

\begin{definition}[Lindblad generator]\label{def:lindblad}
The \emph{Lindblad master equation} for the reduced register state
$\rho_Q = \operatorname{Tr}_{BE}[\rho_{\mathrm{tot}}]$ is
\begin{equation}\label{eq:lindblad}
  \dot{\rho}_Q
  = \mathcal{L}[\rho_Q]
  := -\frac{i}{\hbar}[H_S, \rho_Q]
     + \sum_k \gamma_k \!\left(
         L_k \rho_Q L_k^\dagger
         - \tfrac{1}{2}\{L_k^\dagger L_k, \rho_Q\}
       \right),
\end{equation}
where $H_S \in \mathcal{B}(\mathcal{H}_Q)$ is the system Hamiltonian,
$L_k$ are Lindblad operators, and $\gamma_k \geq 0$ are decay rates.
\end{definition}

\begin{definition}[Hybrid state and control system]\label{def:hybrid}
The \emph{hybrid state} is the pair $(x, \rho_Q)$ with
$x \in \mathbb{R}^m$ a coarse-grained biological/neural variable and
$\rho_Q \in \mathcal{D}(\mathcal{H}_Q)$ the quantum register state.
Its evolution is governed by
\begin{align}
  \dot{x} &= f(x, u, t) + \xi(t), \label{eq:classical-ode}\\
  \dot{\rho}_Q &= \mathcal{L}[\rho_Q] + \mathcal{C}(x, t), \label{eq:quantum-ode}
\end{align}
where $u(t)$ is the control input, $\xi(t)$ is zero-mean white noise
with covariance $Q_\xi$, and
$\mathcal{C}(x, t) : \mathcal{D}(\mathcal{H}_Q) \to \mathcal{D}(\mathcal{H}_Q)$
is the control superoperator induced by $x$.
\end{definition}

\begin{definition}[Biological-to-quantum channel]\label{def:phi}
The \emph{neural-quantum coupling channel} is a family of CPTP maps
\[
  \Phi_t : \mathcal{B}(\mathcal{H}_B) \to \mathcal{B}(\mathcal{H}_Q),
  \qquad
  \rho_B \mapsto \rho_Q^{\mathrm{in}}(t),
\]
characterised by bandwidth $W_\Phi$ (Hz), latency $\tau_\Phi$ (s),
and depolarising fidelity $\mathcal{F}_\Phi$.
\end{definition}

\begin{definition}[Biological coupling update]\label{def:bio-update}
At each discrete cycle step the quantum register state is updated by
the convex mixture
\begin{equation}\label{eq:bio-update}
  \rho_Q \;\leftarrow\; (1-\lambda)\,\rho_Q + \lambda\,\Phi_t[\rho_B],
  \qquad \lambda \in [0,1],
\end{equation}
where $\lambda$ is the \emph{coupling strength} and
$\Phi_t[\rho_B] \in \mathcal{D}(\mathcal{H}_Q)$ is the quantum state
injected by the biological channel.  The map is CPTP for every
$\lambda \in [0,1]$ because it is a convex combination of two CPTP
maps ($\operatorname{id}$ and $\Phi_t$).
\end{definition}

\begin{definition}[Fractal code family]\label{def:fractal-code}
A \emph{fractal quantum code family} $\{\mathcal{E}_n\}_{n \geq 0}$ is
defined by the recursive map
\begin{equation}\label{eq:fractal-recursion}
  \mathcal{E}_{n+1} = \mathcal{R}(\mathcal{E}_n \otimes \mathcal{E}_n),
\end{equation}
where $\mathcal{R}$ is a fixed CPTP \emph{concatenation
superoperator}.  The \emph{physical error rate} is $p_{\mathrm{phys}}$,
the \emph{threshold} is $p_* > 0$, and the \emph{logical error rate} at
level $n$ obeys
\begin{equation}\label{eq:pL-scaling}
  p_L(n) \sim A \!\left(\frac{p_{\mathrm{phys}}}{p_*}\right)^{\!\alpha^n},
  \qquad A, \alpha > 0.
\end{equation}
\end{definition}

% ============================================================
\section{System Model}
\label{sec:ch21-model}

\subsection{Total Hamiltonian decomposition}

We decompose
\begin{equation}\label{eq:H-total}
  H_{\mathrm{tot}}
  = H_Q \otimes \mathbf{1}_B \otimes \mathbf{1}_E
  + \mathbf{1}_Q \otimes H_B \otimes \mathbf{1}_E
  + \mathbf{1}_Q \otimes \mathbf{1}_B \otimes H_E
  + H_{QB} + H_{QE} + H_{BE},
\end{equation}
where $H_{QB}$ encodes the neural-quantum coupling and $H_{QE}$ drives
decoherence.  Under the Born--Markov approximation (environment
correlation time $\tau_E \ll T_{\mathrm{cycle}}$), integrating out
$\mathcal{H}_B \otimes \mathcal{H}_E$ yields~\eqref{eq:lindblad} with
additional control injection~\eqref{eq:quantum-ode}.

\subsection{Cycle constraint}

Viable operation requires
\begin{equation}\label{eq:cycle-constraint}
  T_{\mathrm{cycle}} < 300\,\mu\text{s},
  \qquad
  T_2 \gg T_{\mathrm{cycle}},
\end{equation}
where $T_2$ is the transverse coherence time of the physical qubits.
The cycle budget is allocated as in Table~\ref{tab:timing}.

\begin{table}[ht]
\centering
\caption{CIIR controller timing budget ($T_{\mathrm{cycle}} = 300\,\mu$s total).}
\label{tab:timing}
\begin{tabular}{lll}
\toprule
\textbf{Operation} & \textbf{Symbol} & \textbf{Budget} \\
\midrule
Neural signal acquisition & $\tau_{\mathrm{acq}}$ & $\leq 50\,\mu$s \\
BCI channel $\Phi_t$       & $\tau_{\Phi}$         & $\leq 30\,\mu$s \\
Hamiltonian pulse          & $\tau_{\mathrm{pulse}}$ & $\leq 100\,\mu$s \\
Error correction (level $n$) & $\tau_{\mathrm{EC}}(n)$ & $\leq 80\,\mu$s \\
Feedback computation       & $\tau_{\mathrm{FB}}$   & $\leq 20\,\mu$s \\
Readout                    & $\tau_{\mathrm{RO}}$   & $\leq 20\,\mu$s \\
\midrule
Total                      & $T_{\mathrm{cycle}}$   & $300\,\mu$s \\
\bottomrule
\end{tabular}
\end{table}

% ============================================================
\section{Multi-Qubit Entangled Dynamics}
\label{sec:ch21-entangled}

\subsection{Entanglement Hamiltonian}

For an $N$-qubit register with on-site energies $\omega_i$ and
Ising-type couplings $J_{ij}$, the free Hamiltonian is
\begin{equation}\label{eq:H-free}
  H_0 = \sum_{i=1}^{N} \frac{\omega_i}{2}\,\sigma_z^{(i)}
        + \sum_{i < j} J_{ij}\,\sigma_z^{(i)}\sigma_z^{(j)},
\end{equation}
where $\sigma_z^{(i)}$ acts as the Pauli-$Z$ matrix on qubit $i$ and as
the identity on all others.

\subsection{Control Hamiltonians}

Single-qubit drives and two-qubit entangling gates are generated by
\begin{equation}\label{eq:H-ctrl}
  H_{\mathrm{ctrl}}(t) = \sum_{i=1}^{N}
    \Bigl[
      \Omega_i^x(t)\,\sigma_x^{(i)} + \Omega_i^y(t)\,\sigma_y^{(i)}
    \Bigr]
  + \sum_{i < j} g_{ij}(t)\,\bigl(\sigma_x^{(i)}\sigma_x^{(j)}
                                   + \sigma_y^{(i)}\sigma_y^{(j)}\bigr),
\end{equation}
where $\Omega_i^\mu(t)$ are real pulse amplitudes and $g_{ij}(t)$ are
switched exchange couplings (photonic or phononic mediators).

\subsection{Multi-qubit Lindblad equation}

The full open-system dynamics reads
\begin{equation}\label{eq:multi-lindblad}
  \dot{\rho}_Q
  = -\frac{i}{\hbar}\bigl[H_0 + H_{\mathrm{ctrl}}(t),\,\rho_Q\bigr]
  + \sum_{i=1}^{N}\sum_{k} \gamma_k^{(i)}
      \!\left(
        L_k^{(i)} \rho_Q \bigl(L_k^{(i)}\bigr)^\dagger
        - \tfrac{1}{2}
          \bigl\{\bigl(L_k^{(i)}\bigr)^\dagger L_k^{(i)},\, \rho_Q\bigr\}
      \right),
\end{equation}
where typical noise channels include amplitude damping
($L_{\downarrow}^{(i)} = \sqrt{\gamma_1^{(i)}}\,\sigma_-^{(i)}$),
pure dephasing ($L_{\phi}^{(i)} = \sqrt{\gamma_\phi^{(i)}/2}\,\sigma_z^{(i)}$),
and collective dephasing from low-frequency electromagnetic noise.

\subsection{Validity of the Markovian approximation}

The Markovian approximation is justified when
\begin{equation}\label{eq:markov-cond}
  \max_k \gamma_k \cdot T_{\mathrm{cycle}} \ll 1,
\end{equation}
i.e., the decay is negligible on the cycle timescale.  Non-Markovian
contributions dominate when $\tau_E \gtrsim T_{\mathrm{cycle}}$, in
which case~\eqref{eq:multi-lindblad} must be replaced by a
time-convolutionless or Nakajima--Zwanzig master equation.

% ============================================================
\section{Fractal Encoding and Recursive Error Correction}
\label{sec:ch21-fractal}

\subsection{Renormalisation group interpretation}

Define the RG-like map $\mathcal{T}$ on the space of codes by
$\mathcal{T}[\mathcal{E}] = \mathcal{R}(\mathcal{E} \otimes \mathcal{E})$.
The sequence $\mathcal{E}_n = \mathcal{T}^n[\mathcal{E}_0]$ is the
orbit of $\mathcal{E}_0$ under repeated application of $\mathcal{T}$.
Fixed points $\mathcal{E}^* = \mathcal{T}[\mathcal{E}^*]$ correspond to
self-similar codes.

\subsection{Error scaling derivation}

Suppose $\mathcal{E}_0$ corrects all errors of weight $\leq t$ on
$m_0$ physical qubits, with logical error rate $p_1$.  After one
concatenation, the $[[m_0^2, 1, d_0^2]]$ code corrects errors of weight
$\leq 2t+1$ and logical error rate satisfies
\begin{equation}\label{eq:pL-step}
  p_{L,n+1} \leq c\, p_{L,n}^2
  \quad\Longleftrightarrow\quad
  \frac{p_{L,n+1}}{p_*} \leq \left(\frac{p_{L,n}}{p_*}\right)^2,
  \quad c = \frac{1}{p_*}.
\end{equation}
Iterating from $p_{L,0} = p_{\mathrm{phys}}$:
\begin{equation}\label{eq:pL-iterate}
  p_L(n) = p_* \left(\frac{p_{\mathrm{phys}}}{p_*}\right)^{2^n}.
\end{equation}
The general family~\eqref{eq:pL-scaling} follows with $\alpha = 2$ and
$A = p_*$ for standard concatenated codes, or $\alpha > 2$ for
higher-rate fractal constructions in which $\mathcal{R}$ concatenates
more than two copies.

\subsection{Code distance and overhead}

At level $n$, the code encodes $k = 1$ logical qubit into
$m_0^{2^n}$ physical qubits with code distance $d_n = d_0^{2^n}$.
Hence the overhead scales as $\mathcal{O}(m_0^{2^n})$, which is
super-polynomial in $n$ but exponentially smaller than arbitrary
unencoded computation at the same fidelity.

% ============================================================
\section{CIIR Controllability Theorem}
\label{sec:ch21-controllability}

\begin{theorem}[CIIR Controllability]\label{thm:controllability}
Let $\{H_0, H_{\mathrm{ctrl}}^{(1)},\dots,H_{\mathrm{ctrl}}^{(K)}\}
\subset \mathfrak{u}(2^N)$ be the set of generators available to the
CIIR controller.  The quantum register is \emph{universally
controllable}---i.e., any unitary $U \in \mathrm{SU}(2^N)$ is
reachable from the identity by the controlled Schr\"odinger flow---if
and only if the dynamical Lie algebra satisfies
\begin{equation}\label{eq:lie-closure}
  \mathfrak{g}
  := \operatorname{Lie}\!\bigl\{
       -\tfrac{i}{\hbar}H_0,\,
       -\tfrac{i}{\hbar}H_{\mathrm{ctrl}}^{(k)},\,
       k = 1,\dots,K
     \bigr\}
  = \mathfrak{su}(2^N).
\end{equation}
\end{theorem}

\begin{proof}
The result is a direct application of the Jurdjevic--Sussmann theorem
for bilinear control systems on compact Lie groups~\cite{JS1972}.
Specifically, consider the controlled Schr\"odinger equation
\begin{equation}\label{eq:schrodinger-ctrl}
  \dot{U}(t) = -\frac{i}{\hbar}
    \Bigl[H_0 + \textstyle\sum_k u_k(t)\,H_{\mathrm{ctrl}}^{(k)}\Bigr]
    U(t),
  \quad U(0) = \mathbf{1}.
\end{equation}
This is a right-invariant control-affine system on the compact Lie group
$G = \mathrm{SU}(2^N)$ with Lie algebra $\mathfrak{g}_0 = \mathfrak{su}(2^N)$.

\medskip\noindent
($\Rightarrow$) If $\mathfrak{g} = \mathfrak{su}(2^N)$, the
Jurdjevic--Sussmann theorem guarantees that the reachable set from the
identity equals $\mathrm{SU}(2^N)$.

\medskip\noindent
($\Leftarrow$) If $\mathfrak{g} \subsetneq \mathfrak{su}(2^N)$, the
orbit of the identity is a proper connected subgroup of $\mathrm{SU}(2^N)$,
so some unitaries are unreachable regardless of the control signals
$u_k(t)$.

\medskip\noindent
In practice, $\mathfrak{g} = \mathfrak{su}(2^N)$ is verified
constructively: $H_0$ contains $ZZ$ couplings that generate the
Cartan subalgebra; the single-qubit drives $\sigma_{x,y}^{(i)}$ and
their iterated Lie brackets
$[[\sigma_x^{(i)}, \sigma_z^{(i)}], H_0] \ni \sigma_y^{(j)}$, together
with the exchange terms $g_{ij}\sigma_x^{(i)}\sigma_x^{(j)}$, span all
$\mathfrak{su}(2) \hookrightarrow \mathfrak{su}(2^N)$ factors.  One
verifies dimension: $\dim \mathfrak{su}(2^N) = 4^N - 1$, and the
generated algebra equals this dimension by explicit bracket closure,
completing the proof.
\end{proof}

\begin{corollary}
Under the conditions of Theorem~\ref{thm:controllability}, any CPTP
target map $\mathcal{T}_{\mathrm{target}}$ that can be expressed as a
convex mixture of unitary channels is reachable by the CIIR controller
via time-optimal pulse sequences.
\end{corollary}

% ============================================================
\section{CIIR Threshold Theorem}
\label{sec:ch21-threshold}

\begin{theorem}[CIIR Threshold]\label{thm:threshold}
For the fractal code family defined in Definition~\ref{def:fractal-code}
with recursion~\eqref{eq:fractal-recursion}, there exists a threshold
$p_* \in (0, 1)$ such that:
\begin{enumerate}[label=(\roman*), nosep]
  \item If $p_{\mathrm{phys}} < p_*$, then
        $p_L(n) \to 0$ super-exponentially as $n \to \infty$.
  \item If $p_{\mathrm{phys}} \geq p_*$, then the logical error rate
        does not decrease with code level.
\end{enumerate}
Moreover, for $p_{\mathrm{phys}} < p_*$:
\begin{equation}\label{eq:pL-bound}
  p_L(n) \leq A\!\left(\frac{p_{\mathrm{phys}}}{p_*}\right)^{\!\alpha^n},
  \qquad A = p_*, \quad \alpha \geq 2.
\end{equation}
\end{theorem}

\begin{proof}
\textbf{Step 1: Existence of $p_*$.}
The concatenation map $\mathcal{T}$ acts on error rates as
$f: [0,1] \to [0,1]$, $f(p) = c\, p^2$ where $c > 0$ depends on
the specific code and noise model (for independent depolarising noise
on a stabiliser code with threshold $p_* = c^{-1}$, this is exact).
The fixed points of $f$ are $p = 0$ (trivially stable) and $p = 1/c$
(unstable, hence $p_* = 1/c$).

\medskip\noindent
\textbf{Step 2: Below threshold.}
For $p = p_{\mathrm{phys}} < p_*$, let $r = p/p_* < 1$.  Then
$f(p)/p_* = c(p_*/c\cdot r)^2 / p_* = r^2$, so $p_L(n)/p_* =
r^{2^n} \to 0$.  By~\eqref{eq:pL-iterate} with $A = p_*$, $\alpha = 2$.

\medskip\noindent
\textbf{Step 3: Above threshold.}
For $p \geq p_*$, $f(p) \geq p$ (since $f(p) - p = c p^2 - p = p(cp -
1) \geq 0$ for $p \geq 1/c = p_*$), so the error rate is non-decreasing
under concatenation.

\medskip\noindent
\textbf{Step 4: Super-exponential decay ($\alpha > 2$).}
For codes in which $\mathcal{R}$ concatenates $q \geq 3$ copies of
$\mathcal{E}_n$, the error map satisfies $f(p) \leq c' p^q$, giving
$p_L(n) \leq A (p/p_*)^{q^n}$, so $\alpha = q > 2$.

\medskip\noindent
\textbf{Step 5: Physical realisability with biological noise.}
Biological noise contributions from thermal fluctuations, vibrational
modes, and electromagnetic interference enter as additional Lindblad
channels with effective rates $\gamma_k^{\mathrm{bio}}$.  These raise
$p_{\mathrm{phys}}$ by
$\delta p \sim \sum_k \gamma_k^{\mathrm{bio}} \cdot T_{\mathrm{cycle}}$.
The condition $\gamma_k T_{\mathrm{cycle}} \ll 1$ (Axiom A5,
Eq.~\eqref{eq:markov-cond}) guarantees $\delta p \ll p_*$, so the
effective physical error rate remains below threshold.
\end{proof}

\begin{remark}
Theorem~\ref{thm:threshold} requires that noise channels acting on
different physical qubits be \emph{independent} at the code block
level.  Correlated biological noise---e.g., low-frequency vibrational
bath modes coupling many qubits simultaneously---can violate this
assumption and raise the effective threshold.  Handling correlated noise
is an open problem; see Section~\ref{sec:ch21-failure}.
\end{remark}

% ============================================================
\section{CIIR Phase Theorem}
\label{sec:ch21-phase}

The Controllability Theorem (Section~\ref{sec:ch21-controllability}) and
the Threshold Theorem (Section~\ref{sec:ch21-threshold}) each identify a
distinct necessary condition for viable CIIR operation.  The following
master theorem unifies them into a single phase-diagram result.

\begin{theorem}[CIIR Phase Theorem]\label{thm:phase}
Let the physical parameters of the CIIR controller be given by
\[
  \mathbf{q}
  := \bigl(
       \mathfrak{g},\;
       p_{\mathrm{phys}},\;
       \kappa,\;
       \tau_{\mathrm{FB}},\;
       W_\Phi
     \bigr),
\]
where $\mathfrak{g}$ is the dynamical Lie algebra, $p_{\mathrm{phys}}$
the physical error rate, $\kappa > 0$ the Lyapunov decay rate,
$\tau_{\mathrm{FB}}$ the total feedback latency, and $W_\Phi$ the BCI
bandwidth.  Define the \emph{CIIR-stable phase} to be the subset
\begin{equation}\label{eq:stable-phase}
  \mathcal{S}
  := \bigl\{
       \mathbf{q}
       \;\big|\;
       \mathfrak{g} = \mathfrak{su}(2^N),\;
       p_{\mathrm{phys}} < p_*,\;
       \tau_{\mathrm{FB}} < T_{\mathrm{cycle}}/2,\;
       W_\Phi \geq \kappa/\pi
     \bigr\}.
\end{equation}
Then:
\begin{enumerate}[label=(\roman*), nosep]
  \item \textbf{(Existence of stable phase.)} $\mathcal{S}$ is non-empty:
        the conjunction of the four conditions is simultaneously satisfiable
        for physically realisable parameters.
  \item \textbf{(Universality.)} For any $\mathbf{q} \in \mathcal{S}$, the
        CIIR controller achieves universal unitary control, super-exponential
        logical error suppression, and exponential convergence to $(x^*,\rho^*)$.
  \item \textbf{(Phase boundary.)} The complement $\mathcal{S}^c$ is the
        union of four failure phases FM1--FM4 (Section~\ref{sec:ch21-failure}),
        which are disjoint on the parameter boundary and generically co-dimension-one
        hypersurfaces in $\mathbf{q}$-space.
  \item \textbf{(Critical scaling.)}  On the boundary $p_{\mathrm{phys}} \to p_*^-$,
        the logical error rate at level $n$ satisfies
        \[
          p_L(n) \sim p_* \!\left(\frac{p_{\mathrm{phys}}}{p_*}\right)^{\!\alpha^n},
        \]
        and the ``fault-tolerance correlation length''
        $\xi := |\log p_L(n)|^{1/n} \to \alpha$ as $n \to \infty$.
\end{enumerate}
\end{theorem}

\begin{proof}
\textbf{(i) Existence.}  Take $N = 5$, independent depolarising noise
with $p_{\mathrm{phys}} = 10^{-3}$, $p_* = 10^{-2}$ (standard stabiliser
threshold), $\tau_{\mathrm{FB}} = 50\,\mu\text{s}$,
$T_{\mathrm{cycle}} = 300\,\mu\text{s}$, $\kappa = 2\pi \times 1\,\text{kHz}$,
and $W_\Phi = 4\,\text{kHz}$.  Direct substitution verifies all four
inequalities in~\eqref{eq:stable-phase}.

\medskip\noindent
\textbf{(ii) Universality.}  By Theorem~\ref{thm:controllability},
$\mathfrak{g} = \mathfrak{su}(2^N)$ implies any target unitary is
reachable.  By Theorem~\ref{thm:threshold}, $p_{\mathrm{phys}} < p_*$
implies $p_L(n) \to 0$ super-exponentially.  By
Proposition~\ref{prop:lyapunov}, $\kappa > 0$ and $W_\Phi \geq \kappa/\pi$
(Nyquist condition, \eqref{eq:min-bandwidth}) imply exponential convergence
to $(x^*, \rho^*)$.  The latency condition $\tau_{\mathrm{FB}} < T_{\mathrm{cycle}}/2$
prevents FM2 instability, so the feedback loop remains causal.

\medskip\noindent
\textbf{(iii) Phase boundary.}  Each failure condition in FM1--FM4 defines
a boundary hypersurface:
\[
  \partial\mathcal{S}
  = \bigl\{\max_k \gamma_k T_{\mathrm{cycle}} = 1\bigr\}
  \cup \bigl\{p_{\mathrm{phys}} = p_*\bigr\}
  \cup \bigl\{\tau_{\mathrm{FB}} = T_{\mathrm{cycle}}/2\bigr\}
  \cup \bigl\{W_\Phi = \kappa/\pi\bigr\}.
\]
Each surface is co-dimension one in the four-dimensional parameter space
$\mathbf{q} \in \mathbb{R}_{\geq 0}^4$.  The four conditions are
generically independent (their Jacobians are linearly independent at a
generic boundary point), so FM1--FM4 boundaries intersect only in
measure-zero subsets.

\medskip\noindent
\textbf{(iv) Critical scaling.}  Follows directly from
Theorem~\ref{thm:threshold} with $A = p_*$ and $\alpha \geq 2$.  The
correlation length $\xi$ is defined by analogy with the RG interpretation
(Section~\ref{sec:ch21-discussion}).
\end{proof}

\begin{remark}
Theorem~\ref{thm:phase} is the CIIR analogue of a quantum phase transition
diagram: inside $\mathcal{S}$ the controller is in the ``fault-tolerant
and controllable'' phase; outside $\mathcal{S}$ it resides in one of four
distinct error phases.  The boundaries are sharp in the thermodynamic
limit $N \to \infty$ and are smeared to crossovers for finite $N$.
\end{remark}

% ============================================================
\section{Stability and Attractor Analysis}
\label{sec:ch21-stability}

\subsection{Attractors of the Lindbladian flow}

The semigroup $\{e^{t\mathcal{L}}\}_{t \geq 0}$ generated by the
Lindblad superoperator $\mathcal{L}$ (Eq.~\eqref{eq:lindblad}) is a
contractive completely positive map for each $t \geq 0$.  By the
Perron--Frobenius theorem for positive maps, the peripheral spectrum
$\sigma_{\mathrm{per}}(\mathcal{L}) = \{\lambda \in \sigma(\mathcal{L}) :
\operatorname{Re}\lambda = 0\}$ determines the set of invariant states.
If $\sigma_{\mathrm{per}}(\mathcal{L}) = \{0\}$ (unique invariant state),
the flow is \emph{globally attractive} to the unique steady state
$\rho_Q^*$; otherwise multiple invariant subspaces coexist.

\subsection{Fixed-point condition}

Define the composed map $\Psi_t := \Phi_t \circ e^{t\mathcal{L}}$
(apply Lindblad flow then inject biological input).  The CIIR
\emph{operating fixed point} $\rho^*$ satisfies
\begin{equation}\label{eq:fixed-point}
  \Psi_t(\rho^*) = \rho^*, \qquad \forall\, t.
\end{equation}
By Schauder's fixed-point theorem applied to the compact convex set
$\mathcal{D}(\mathcal{H}_Q)$, at least one such fixed point exists for
every continuous $\Psi_t$ (Lemma~3 of Chapter~\ref{ch:multi-qubit-controller}
restates this).

\subsection{Lyapunov stability of the hybrid system}

\begin{proposition}[Lyapunov stability]\label{prop:lyapunov}
Suppose there exists a smooth Lyapunov functional
\begin{equation}\label{eq:lyapunov}
  V(x, \rho_Q)
  = \tfrac{1}{2}\|x - x^*\|_P^2
  + D_{\mathrm{re}}(\rho_Q\|\rho^*),
\end{equation}
where $\|\cdot\|_P^2 = (\cdot)^\top P(\cdot)$ with $P > 0$, and
$D_{\mathrm{re}}(\rho\|\sigma) = \operatorname{Tr}[\rho(\log\rho - \log\sigma)]$
is the quantum relative entropy.  If the controlled system satisfies
\begin{equation}\label{eq:lyapunov-decay}
  \frac{d}{dt} V(x(t), \rho_Q(t))
  \leq -\kappa V(x(t), \rho_Q(t)),
  \quad \kappa > 0,
\end{equation}
then $(x^*, \rho^*)$ is exponentially stable with rate $\kappa$.
\end{proposition}

\begin{proof}
Differentiating $V$ along trajectories of~\eqref{eq:classical-ode}--\eqref{eq:quantum-ode}:
\begin{align}
  \dot{V}
  &= (x - x^*)^\top P \dot{x}
  + \frac{d}{dt} D_{\mathrm{re}}(\rho_Q\|\rho^*) \notag\\
  &= (x - x^*)^\top P [f(x, u, t) + \xi]
  + \operatorname{Tr}\bigl[
      (\log\rho_Q - \log\rho^*)\,
      (\mathcal{L}[\rho_Q] + \mathcal{C}(x,t))
    \bigr].
\end{align}
Choose the feedback $u = u^*(x, \rho_Q)$ to satisfy
$(x - x^*)^\top P f(x, u^*, t) \leq -\lambda_{\min}(P) \|x - x^*\|^2$
(linear stabilising feedback for the $x$-subsystem).  For the quantum
part, the contractivity of $e^{t\mathcal{L}}$ in relative entropy
\[
  D_{\mathrm{re}}\!\bigl(e^{t\mathcal{L}}[\rho]\,\big\|\,\rho^*\bigr)
  \leq D_{\mathrm{re}}(\rho\|\rho^*)
\]
ensures the quantum contribution is non-increasing; the control
injection $\mathcal{C}(x,t)$ is designed to make it strictly
decreasing.  Under these conditions~\eqref{eq:lyapunov-decay} holds,
and Barbalat's lemma implies exponential convergence to $(x^*, \rho^*)$.
\end{proof}

% ============================================================
\section{Control Theory Formulation}
\label{sec:ch21-control}

\subsection{Optimal control problem}

Given a target unitary $U_{\mathrm{target}} \in \mathrm{SU}(2^N)$ and a
time horizon $[0, T]$ with $T \leq T_{\mathrm{cycle}}$, define the
cost functional
\begin{equation}\label{eq:cost}
  J[\mathbf{u}]
  = \frac{1}{2}
    \bigl\|U(T) - U_{\mathrm{target}}\bigr\|_F^2
  + \frac{\mu}{2} \int_0^T \|\mathbf{u}(t)\|^2\,dt,
\end{equation}
where $\|\cdot\|_F$ is the Frobenius norm, $\mathbf{u}(t) =
(\Omega_1^x, \Omega_1^y, \dots, \Omega_N^y, g_{ij}, \dots)$ is the
concatenated control vector, and $\mu > 0$ penalises power.

Necessary conditions for optimality (Pontryagin Maximum Principle)
yield the \emph{GRAPE equations}:
\begin{equation}\label{eq:grape}
  \frac{\partial J}{\partial \mathbf{u}(t)}
  = -\frac{i}{\hbar}
    \operatorname{Tr}\!\bigl[
      \Lambda^\dagger(t)\,
      \frac{\partial H_{\mathrm{ctrl}}}{\partial \mathbf{u}}\,
      U(t)
    \bigr]
    + \mu\,\mathbf{u}(t) = 0,
\end{equation}
where $\Lambda(t)$ satisfies the co-state equation
$\dot{\Lambda} = -\frac{i}{\hbar}[H_0 + H_{\mathrm{ctrl}}(t),\Lambda]$
with $\Lambda(T) = -(U(T) - U_{\mathrm{target}})$.

\subsection{Real-time adaptive feedback}

The feedback controller observes the syndrome measurements from the
fractal code and the classical state $x$ from the BCI, and updates the
control pulses via
\begin{equation}\label{eq:feedback-law}
  \mathbf{u}(t) = \mathbf{u}^*(t) + K\!\bigl[\hat{x}(t) - x^*\bigr]
                 + K_Q\!\bigl[\mathcal{M}[\rho_Q(t)] - \mathcal{M}[\rho^*]\bigr],
\end{equation}
where $\mathbf{u}^*(t)$ is the pre-computed GRAPE solution, $K$ and
$K_Q$ are gain matrices designed via the Lyapunov method of
Proposition~\ref{prop:lyapunov}, and $\mathcal{M}$ denotes the syndrome
measurement map.

% ============================================================
\section{Full CIIR Controller Architecture}
\label{sec:ch21-architecture}

The five-layer architecture is defined as follows.

\paragraph{Layer 1: Molecular qubits (protein-based).}
$N$ protein qubits embedded via $V$ (Definition~\ref{def:protein-qubit})
with individual control by microwave/optical drives.  Physical parameters:
$T_1 \sim 1\,\text{ms}$, $T_2 \sim 100\,\mu\text{s}$,
single-qubit gate fidelity $\mathcal{F}_1 \geq 0.999$.

\paragraph{Layer 2: Coupling network ($J_{ij}$, photonic/phononic modes).}
All-to-all effective couplings mediated by a cavity QED or phononic
crystal bus.  The coupling matrix $J_{ij}$ is programmed by the
feedback controller; two-qubit gate fidelity $\mathcal{F}_2 \geq 0.99$.

\paragraph{Layer 3: Control interface (BCI).}
The neural-quantum channel $\Phi_t$ (Definition~\ref{def:phi})
translates decoded neural intent signals $x$ into control
Hamiltonians $H_{\mathrm{ctrl}}(t)$.  Required: bandwidth
$W_\Phi \geq 1/T_{\mathrm{cycle}} \approx 3.3\,\text{kHz}$, latency
$\tau_\Phi \leq 30\,\mu\text{s}$.

\paragraph{Layer 4: Error correction (fractal code $\mathcal{E}_n$).}
Level-$n$ fractal code (Definition~\ref{def:fractal-code}) with
threshold $p_* \sim 10^{-2}$ (stabiliser code family).  Syndrome
extraction and correction sub-cycles run within the $80\,\mu$s budget.

\paragraph{Layer 5: Feedback controller (real-time adaptive).}
GRAPE optimal pulses~\eqref{eq:grape} with real-time adaptive
corrections~\eqref{eq:feedback-law}.  Hardware: FPGA with $\leq 20\,\mu$s
classical computation latency.

\subsection{Signal flow}

The formal signal flow is:
\begin{equation}\label{eq:signal-flow}
  \underbrace{x(t)}_{\text{neural state}}
  \xrightarrow{\Phi_t}
  \underbrace{\rho_Q^{\mathrm{in}}(t)}_{\text{injected state}}
  \xrightarrow{e^{\Delta t\,\mathcal{L}}}
  \underbrace{\rho_Q(t+\Delta t)}_{\text{evolved state}}
  \xrightarrow{\mathcal{E}_n}
  \underbrace{\bar{\rho}_Q(t+\Delta t)}_{\text{error-corrected state}}
  \xrightarrow{\mathcal{M}}
  \underbrace{s(t)}_{\text{syndrome}}
  \xrightarrow{K,K_Q}
  \underbrace{\mathbf{u}(t+\Delta t)}_{\text{updated control}}.
\end{equation}

\subsection{Minimum viable bandwidth}

For stable feedback control the Nyquist condition requires
$W_\Phi \geq 2 B_{\mathrm{ctrl}}$, where $B_{\mathrm{ctrl}}$ is the
control bandwidth set by the fastest desired state trajectory.  Setting
$B_{\mathrm{ctrl}} = \kappa/2\pi$ (from the Lyapunov decay rate), the
minimum viable bandwidth is
\begin{equation}\label{eq:min-bandwidth}
  W_\Phi \geq \frac{\kappa}{\pi}.
\end{equation}
For $\kappa = 2\pi \times 1\,\text{kHz}$ this gives $W_\Phi \geq 2\,\text{kHz}$,
comfortably within the allocated $\tau_\Phi = 30\,\mu$s budget.

% ============================================================
\section{Timing, Bandwidth, and Physical Constraints}
\label{sec:ch21-constraints}

\subsection{Decoherence parameter}

Viable operation requires (from~\eqref{eq:markov-cond}):
\begin{equation}\label{eq:decoherence-param}
  \epsilon_k := \gamma_k T_{\mathrm{cycle}} \ll 1,
  \quad \forall k.
\end{equation}
For $T_{\mathrm{cycle}} = 300\,\mu$s and target $\epsilon_k < 10^{-2}$,
the required decay rates satisfy $\gamma_k < 33\,\text{s}^{-1}$, corresponding
to $T_1, T_2 > 30\,\text{ms}$---achievable in cryogenic protein-qubit
platforms (see Section~\ref{sec:ch21-experiment}).

\subsection{Channel capacity}

The quantum channel capacity of $\Phi_t$ under a depolarising noise
model with fidelity $\mathcal{F}_\Phi$ and bandwidth $W_\Phi$ is
bounded by
\begin{equation}\label{eq:channel-capacity}
  C_Q(\Phi_t)
  \leq W_\Phi \cdot h_q(\mathcal{F}_\Phi),
\end{equation}
where $h_q(\mathcal{F}) = \log_2 2 - h_2\!\left(\frac{1-\mathcal{F}}{3}\right)
- \frac{1-\mathcal{F}}{3}\log_2 3$ is the hashing bound for the
depolarising channel and $h_2(p) = -p\log_2 p - (1-p)\log_2(1-p)$.
The mutual information per cycle is
\begin{equation}\label{eq:mutual-info}
  I(B \to Q)_{\mathrm{cycle}}
  \leq C_Q(\Phi_t) \cdot T_{\mathrm{cycle}}.
\end{equation}

% ============================================================
\section{Failure Modes and Phase Boundaries}
\label{sec:ch21-failure}

We enumerate four critical failure modes as strict inequality conditions.

\begin{description}
  \item[FM1: Decoherence-dominated regime.]
        Failure when $\max_k \gamma_k T_{\mathrm{cycle}} \geq 1$,
        equivalently $T_2 \leq T_{\mathrm{cycle}}$.  In this regime the
        state completely decoheres within one cycle and no fault-tolerant
        operation is possible.

  \item[FM2: Control latency instability.]
        Failure when the total feedback latency
        $\tau_{\mathrm{FB}} = \tau_\Phi + \tau_{\mathrm{comp}} \geq T_{\mathrm{cycle}}/2$.
        The Nyquist condition is violated; the feedback loop destabilises
        rather than corrects.

  \item[FM3: Error correction threshold failure.]
        Failure when $p_{\mathrm{phys}} \geq p_*$.  By
        Theorem~\ref{thm:threshold}(ii), concatenation increases rather
        than decreases the logical error rate.  Equivalently, the
        renormalisation-group flow is directed away from the zero
        fixed point.

  \item[FM4: Information bottleneck collapse.]
        Failure when $I(B \to Q)_{\mathrm{cycle}} < I_{\mathrm{min}}$,
        where $I_{\mathrm{min}}$ is the minimum mutual information
        required for stable feedback (set by the entropy production rate
        of the target task).  This occurs when
        $W_\Phi < \kappa/\pi$ (Eq.~\eqref{eq:min-bandwidth}) or
        $\mathcal{F}_\Phi$ drops below the quantum capacity threshold.
\end{description}

\noindent
The four conditions partition the parameter space
$(\gamma_{\max}, \tau_{\mathrm{FB}}, p_{\mathrm{phys}}, W_\Phi)$
into an operational region (none of FM1--FM4 satisfied) and failure
phases.

% ============================================================
\section{Experimental Realization Pathway}
\label{sec:ch21-experiment}

\subsection{Stage 1: Single protein qubit validation}

\textbf{Observables:} $T_1$, $T_2$, single-qubit gate fidelity
$\mathcal{F}_1$, readout contrast.

\textbf{Success criteria:} $T_2 > 1\,\text{ms}$, $\mathcal{F}_1 > 0.999$,
readout fidelity $> 0.99$.

\textbf{Falsifiability:} If $T_2 < 100\,\mu\text{s}$ under all
protein-engineering protocols, Axiom~A4 fails and the protein-qubit
platform is ruled out.

\subsection{Stage 2: Two-to-ten qubit coupling demonstration}

\textbf{Observables:} Two-qubit Bell state fidelity, entanglement
negativity, cross-resonance coupling rate $g_{ij}$.

\textbf{Success criteria:} Bell state fidelity $> 0.99$, $g_{ij} >
10\,\text{kHz}$.

\textbf{Falsifiability:} If entanglement fidelity cannot exceed
$0.9$ for any two-qubit geometry, the Lie algebra closure
(Theorem~\ref{thm:controllability}) is not achievable with these
physical couplings.

\subsection{Stage 3: Closed-loop control with simulated neural input}

\textbf{Observables:} Gate fidelity under GRAPE, FPGA feedback latency,
error-correction cycle time.

\textbf{Success criteria:} Gate fidelity $> 0.99$ with real-time
feedback, total cycle $\leq 300\,\mu\text{s}$, FM2 condition
$\tau_{\mathrm{FB}} < 150\,\mu\text{s}$ satisfied.

\textbf{Falsifiability:} If latency consistently exceeds
$T_{\mathrm{cycle}}/2$ regardless of hardware, FM2 condemns
this architecture.

\subsection{Stage 4: Real biological interface integration}

\textbf{Observables:} Mutual information $I(B \to Q)$, closed-loop
state fidelity, BCI-to-quantum latency $\tau_\Phi$.

\textbf{Success criteria:} $I(B \to Q) > I_{\mathrm{min}}$,
$\tau_\Phi \leq 30\,\mu\text{s}$, stable operation for $> 10^3$
consecutive cycles.

\textbf{Falsifiability:} If $I(B \to Q) < I_{\mathrm{min}}$ under
all BCI modalities (EEG, ECoG, spike-sorted single units), the
channel capacity bound~\eqref{eq:channel-capacity} is insufficient
for feedback stabilisation and the hybrid architecture fails.

% ============================================================
\section{Discussion and Implications}
\label{sec:ch21-discussion}

\subsection{Renormalisation group unification}

The fractal code recursion~\eqref{eq:fractal-recursion} is formally
equivalent to a discrete RG transformation on the space of quantum
codes.  The threshold $p_*$ is the unstable fixed point of the map $f$,
and the region $p < p_*$ is the basin of attraction of the fault-tolerant
fixed point $p = 0$.  This places quantum fault tolerance within the
universality framework of critical phenomena: different fractal code
families sharing the same $\alpha$ belong to the same \emph{CIIR
universality class}.

\subsection{Dissipative quantum computing}

The Lindblad dynamics~\eqref{eq:multi-lindblad} can be viewed as
\emph{dissipative quantum computing}: steady states of engineered
Lindblad operators encode the computation result.  The CIIR controller
exploits this by (a) engineering $\mathcal{L}$ so that $\rho^*$ encodes
the target logical state, and (b) using $\mathcal{C}(x,t)$ to steer
the system toward that attractor.

\subsection{Universality class}

We conjecture (and leave as an open problem) that all CIIR systems
satisfying Theorems~\ref{thm:controllability}--\ref{thm:threshold} and
Proposition~\ref{prop:lyapunov} belong to a single universality class
characterised by the fixed-point scaling exponent $\alpha$ and the
Lyapunov decay rate $\kappa$.  Systems in the same class share the same
critical scaling of $p_L(n)$ near threshold and the same divergence of
the correlation length $\xi \sim |p - p_*|^{-\nu}$.

\subsection{Outlook}

The architecture presented here bridges quantum information theory,
control theory, and biophysics through the CIIR categorical language.
Key open questions include:
\begin{itemize}[nosep]
  \item Construction of explicit fractal stabiliser codes with $\alpha > 2$
        that are also space-efficient.
  \item Experimental realisation of protein-qubit $T_2 > 1\,\text{ms}$.
  \item Formal characterisation of the CIIR universality class and its
        relationship to topological phases.
  \item Extension to non-Markovian regimes relevant to in-vivo biological
        interfaces.
\end{itemize}
